\documentclass[../main.tex]{subfiles}

\begin{document}
\section{Casos d'ús}
\textbf{Nom:} Crear perfil\\
\textbf{Actor:} Usuari\\
\textbf{Comportament:}
\begin{enumerate}
    \item L’usuari informa del nom de l’usuari per poder respondre
    \item L’usuari pot decidir en respondre amb un email (registrat) o anònimament
    \item El sistema comprova les dades introduïdes i registra un nou perfil
\end{enumerate}
\textbf{Errors possibles i solucions}
\begin{itemize}
    \item Si ja existeix el perfil amb el nom d’usuari el sistema informarà l’error i li recomanarà un nou nom d’usuari bàsic
    \item En cas que l’email ja existeixi el sistema informarà de l’error i no podrà registrar-se amb el perfil, en cas que sigui anònim activarà l’apartat 1a)
\end{itemize}

\vspace{0.5 cm}

\textbf{Nom:} Respondre Enquesta\\
\textbf{Actor:} Usuari\\
\textbf{Comportament:}
\begin{enumerate}
    \item L’usuari respondrà la pregunta de l’enquesta depenent del tipus de resposta
    \item El sistema validarà la resposta i passarà a la següent pregunta
\end{enumerate}
\textbf{Errors possibles i solucions}
\begin{itemize}
    \item Tot i que és possible deixar valors nuls, si ingressa algun caràcter especial el sistema informarà de l’error
\end{itemize}

%%%%%%%%%%%%%%%%%%%%%%%%%%%%%%%%%%%%%%%%%%%%%%%%%%%%%%%%%%%%%%%%%%%%%%%%%%%%%%%%%%%%%%%%%%
\vspace{0.5 cm}

\textbf{Nom:} Consultar perfil
\textbf{Actor:} Usuari
\textbf{Comportament:}
\begin{enumerate}
    \item L’usuari demana al sistema informació del perfil com usuari o email en cas que sigui un usuari registrat
\end{enumerate}
\textbf{Errors possibles i solucions}
\begin{itemize}
    \item En cas que l’usuari anònim demani informació de l’email, el sistema informarà de l’error
\end{itemize}

\vspace{0.5 cm}

\textbf{Nom:} Exportar enquesta
\textbf{Actor:} Usuari
\textbf{Comportament:}
\begin{enumerate}
    \item L’usuari demana al sistema l’exportació de l’enquesta en el format demanat (pdf o txt)
    \item El sistema generarà un document i l’enviarà a la ruta que estableixi l’usuari
\end{enumerate}
\textbf{Errors possibles i solucions}
\begin{itemize}
    \item L’usuari pot incloure una ruta que no existeixi i el sistema informarà de l’error.
    \item En cas que l’usuari no indiqui la ruta el sistema el guardarà per defecte a Baixades
\end{itemize}

\vspace{0.5 cm}

\textbf{Nom:} Crear enquesta
\textbf{Actor:} Usuari
\textbf{Comportament:}
\begin{enumerate}
    \item L’usuari crearà una enquesta amb un seguit de preguntes que el mateix definirà
    \item L’usuari podrà definir el tipus de pregunta i el format de les seves respostes
    \item El sistema generarà l’enquesta amb les preguntes definides i les guardarà automàticament
\end{enumerate}
\textbf{Errors possibles i solucions}
\begin{itemize}
    \item Si l’usuari tracta de finalitzar l’enquesta sense establir les preguntes que ha definit el sistema informarà de l’error
    \itemSi l’usuari tracta d’utilitzar caràcters especials en les preguntes el sistema informarà de l’error però li permetrà continuar
    \item Si l’usuari tanca l’enquesta sense finalitzar les preguntes s’activa la 1a)
\end{itemize}

\vspace{0.5 cm}

\textbf{Nom:} Modificar enquesta
\textbf{Actor:} Usuari
\textbf{Comportament:}
\begin{enumerate}
    \item L’usuari pot modificar la pregunta de les seves enquestes i en cas que siguin diferents es permetrà el canvi d’alguns tipus de resposta
\end{enumerate}
\textbf{Errors possibles i solucions}
\begin{itemize}
    \item Si l’usuari tracta de modificar el número de preguntes el sistema informarà de l’error i l'informarà sobre la possibilitat de fer una nova enquesta
\end{itemize}

\vspace{0.5 cm}

\textbf{Nom:} Carregar enquesta
\textbf{Actor:} Usuari
\textbf{Comportament:}
\begin{enumerate}
    \item L’usuari podrà carregar alguna enquesta que hagi fet o fins i tot alguna incompleta per terminar-la de fer
\end{enumerate}
\textbf{Errors possibles i solucions}
no hay??

\vspace{0.5 cm}

\textbf{Nom:} Guardar enquesta
\textbf{Actor:} Usuari
\textbf{Comportament:}
\begin{enumerate}
    \item L’usuari li demanarà el sistema que la guardi perquè sigui contestada per altres
    \item L’usuari podrà guardar-la en qualsevol moment inclòs si està incomplet
\end{enumerate}
\textbf{Errors possibles i solucions}
\begin{itemize}
    \item En cas que l’usuari tanqui l’enquesta o hi hagi alguna fallada el sistema la guardarà automàticament
\end{itemize}

\vspace{0.5 cm}

\textbf{Nom:} Sortir
\textbf{Actor:} Usuari
\textbf{Comportament:}
\begin{enumerate}
    \item L’usuari podrà sortir de l’enquesta en qualsevol moment
\end{enumerate}
\textbf{Errors possibles i solucions}
\begin{itemize}
    \item En cas que l’usuari està responent o creant alguna enquesta el sistema el guardara automàticament
\end{itemize}
















\end{document}
